\chapter{El sistema de transporte. Modos}
\section{Definición del Sistema de Transporte}
Acción conjunta de los distintos elementos que lo integran: infraestructura, vehículos, empresas, administración, etc., para dar satisfacción a las necesidades de movilidad de la colectividad, en un espacio y tiempo y tiempo determinados.

\section{Principales Características del Sistema de Transporte}
\begin{enumerate}
    \item Gran dependencia del resto de las actividades económicas, al no ser un fin en sí mismo. Dimensionado.
    \item Es un sector muy capitalizado. Instalaciones usmamente costosas.
    \item Discontinuidad de la infraestructura.
    \item Tratamiento diferenciado de la infraestructura y de los vehículos.
    \item El transporte no puede almacenarse.
    \item Descompensación en determinadas corrientes de tráfico. Movimientos en vacío.
    \item Fluctuaciones temporales de la demanda.
    \item Fiscalidad diferente. Los precios no reflejan los costes reales.
    \item Produce fuertes externalidades: costes y beneficios que no afectan al productor ni al consumidor (Madrid central, Puertollano AVE).
    \item Dificultad para definir la unidad de medida.
    \item Información del sector: escasa, poco fiable y a veces contradictoria (Observatorios).
    \item Imperfección del mercado. Consecuencia de todo lo indicado anteriormente.
\end{enumerate}

\section{Modos de Transporte}

Para cada modo de transporte existe, al menos, una Organización de Transporte Internacional para las partes que intervienen en ese modo. Como ejemplos:
\begin{itemize}
    \item Carretera: IRU (Organización Internacional del Transporte por Carretera).
    \item Ferrocarril: UIC (Unión INternacional de Ferrocarriles).
    \item Aéreo: IATA (Asociación Internacional de Transporte Aéreo).
    \item Marítimo: OMI (Organización Marítima Internacional).
    \item Transporte Mundial de Mercancías: FIATA (Federación Internacional de Asociaciones de Agentes de Carga).
\end{itemize}

\section{Características de la Oferta del Transporte}
\begin{enumerate}
    \item Condiciona el crecimiento de la movilidad.
    \item ``Toda oferta crea su propia demanda''. [Jean-Baptista, SCJ. (1767 - 1830)] EJ.: AVE, VTC.
    \item Cantidad y calidad de los servicios de transporte que pueden prestar los transportistas en un territorio (país) y momento determinado (toneladas ferroviarias, plazas de autocar ó avión, etc.).
    \item La infraestructura puede limitar la operatividad de la oferta instantánea.
    \item La oferta instantánea es independiente del precio a corto plazo.
    \item Existencia de solapamiento. Competencia entre modos de transportes.
    \item Existencia de una sobreoferta si se quiere tener una alta calidad.
    \item ``Producción conjunta'' dividida en viajeros y mercancías. Problemas de imputación de costos.
    \item Competencia entre los transportes privados y los transportes públicos. Empresas de tamaño muy diferente. (Taxi / VTC).
    \item Discontinuidad de los medios de producción del transporte.
\end{enumerate}

\section{Parámetros de la Oferta de Viajeros}
\begin{itemize}
    \item Seguridad (CSEEA).
    \item Regularidad.
    \item Velocidad Comercial Media.
    \item Comodidad (puerta a puerta).
    \item Oportunidad.
    \item Frecuencia (just in time).
    \item Flexibilidad (independencia).
    \item Versatilidad del Material (vehículos).
    \item Congestión del Tráfico.
    \item Accidentalidad.
    \item Ofertar Sevicios Complementarios (logística...).
\end{itemize}

\section{Características de Calidad en el Transporte de Mercancías}
\begin{itemize}
    \item Velocidad.
    \item Amplitud de Red.
    \item Independencia.
    \item Capacidad.
    \item Frecuencia.
    \item Coste.
\end{itemize}

\subsection{Comparativa Según Modos}
\begin{table}[H]
    \begin{tabular}{| c | c |}
        Modo & Característica: Velocidad \\
        
    \end{tabular}
\end{table}

\section{Carretera}
\subsection{Vehículos}

\subsection{Categorías de Vehículos}
Categoría M:
\begin{itemize}
    \item M1 - Turismo: automóvil, distinto de la motocicleta, especialmente concebido y construido para el transporte de personas y con capacidad hasta nueve plazas, incluido el conductor.
    \item M2 - M3 - Autobús: automóvil concebido y construido para el transporte de personas, con capacidad para más de nueve plazas, incluido el conductor. Se incluye en este término el trolebús, es decir, el vehículo conectado a una línea eléctrica y que no circula por raíles.
\end{itemize}

Categoría N:
\begin{itemize}
    \item N2 - N3 - Camión: automóvil concebido y construido para el transporte de cosas. Se excluye de esta definición la motocicleta de tres ruedas.
\end{itemize}

Categoría O:
\begin{itemize}
    \item Remolque: Vehículo concebido y construido para circular arrastraod por un vehículo motor.
    \item Semirremolque: remolque construido para ser acoplado a un automóvil de tal manera que repose parcialmente sobre éste y que una parte sustancial de su peso y de su carga sea soportada por dicho automóvil.
    \item Tractor industrial: automóvil concebido y construido para realizar, principalmente, el arrastre de un semirremolque.
\end{itemize}

Otras Categorías L, G, especiales.

\subsection{Infraestructura}


\section{Ferrocarril}
\subsection{Material Rodante (Locomotoras, Material Remolcado) e Infraestructuras}

\subsection{Categorías de Material Rodante}
\begin{itemize}
    \item Tren: serie de vehículos acoplados unos con otros que, remolcados por una locomotora o un automotor, conducen viajeros o mercancías de un punto a otro de una línea férrea.
    \item Locomotora: vehículo ferroviario con motor destinado exclusivamente a remolcar otros vehículos ferroviarios por la vía férrea; según sea la energía que utilice, puede ser diésel, eléctrica, etc..., y su transmisión puede ser: eléctrica, hidráulica, hidromecánica o mecánica.
    \item Automotor: tren formado por material autopropulsado, cualquiera que sea el número de motores, remolques o elementos por los que esté compuesto.
    \item Unidad de tren: conjunto de vehículos automotores y remolques que forman una composición indivisible preparada para circular sola o para acoplarse a otras composiciones.
    \item Coche: vehículo ferroviario destinado al transporte de viajeros.
    \item Vagón: unidad móvil ferroviaria sin sistema de propulsión propio destinada al transporte de mercancías.
    \item Material motor: conjunto de locomotoras, automotores y otras unidades de tracción.
    \item Material móvil (o remolcado): conjunto de coches, vagones y otros vehículos sin motor.
    \item Composición del tren: conjunto de vagones o unidades que forman un tren.
    \item Composición reversible: composición de tren indeformable que puede funcionar indistintamente en ambos sentidos.
\end{itemize}

% Compare en términos relativos los 5 modos de transporte de mercancías. Comparelos desde el punto de vista de velocidad.

\section{Transporte Aéreo}

\section{Transporte Marítimo}
\subsection{Modernización de Puertos (PDI)}
\subsubsection{Demanda Previsible}
\begin{itemize}
    \item Los cambios estratégicos y tecnológicos del transporte marítimo y la explotación portuaria. 
    \item La sostenibilidad ambiental de los tráficos.
    \item El fomento del cabotaje europeo como alternativa al tráfico terrestre intracomunitario.
    \item La seguridad y calidad en la prestación de los servicios portuarios.
    \item La mejora de los puertos como intercambiadores modales con mejores conexiones viarias y ferroviarias.
    \item La potenciación de los puertos como plataformas logísticas.
    \item La mejora de las relaciones puerto ciudad.
\end{itemize}

\subsubsection{Inversión Pública}
\begin{itemize}
    \item Construcción de áreas abrigadas y accesos marítimos.
    \item Incorporación de suelo mediante adquisición de terrenos o ganándolos al mar.
    \item Construcción de atraques asociados a servicios portuarios no concesionados.
    \item Instalación y manteniemi
\end{itemize}